\documentclass[11pt]{article}
\usepackage{fontspec}
\usepackage[french,english]{babel}
\usepackage{amsmath}
\usepackage{amssymb}
\usepackage{amsthm}
\usepackage{mathrsfs}
\usepackage{tikz-cd}
\usepackage{enumitem}
\usetikzlibrary{babel}

\newcommand{\printedpage}[1]{\marginpar{\footnotesize p.~#1}}
\newcommand{\msp}[1]{\marginpar{\footnotesize [#1]}}

\newcommand{\Hom}{\mathrm{Hom}}
\newcommand{\Fl}{\mathrm{Fl}}
\newcommand{\Ob}{\mathrm{Ob}}
\newcommand{\Cat}{\mathrm{Cat}}
\newcommand{\Ens}{\mathrm{Ens}}
\newcommand{\pr}{\mathrm{pr}}
\newcommand{\Ker}{\mathrm{Ker}}
\newcommand{\Coker}{\mathrm{Coker}}
\newcommand{\mon}{\mathrm{mon}}
\newcommand{\deq}{\mathrel{\overset{\text{def}}{=}}}

\begin{document}

\printedpage{}
% no page number printed (title page)
% head: (none)

\begin{center}
{\LARGE\textsf{THE DERIVATORS}}

\bigskip

\textbf{ALEXANDRE GROTHENDIECK}

\bigskip\bigskip

{\large\textbf{Appendix to Chapter XII}}
\end{center}

\newpage

\printedpage{1}
% head: A. Grothendieck : Dérivateurs, Ch. XII, Édition des Universités de Montpellier II et Paris VII
% foot: Version 20 juin 2003

\msp{page 1}
\textbf{Proposition} (?) (\footnote{Cf. pages 11 ff. for the key counterexamples.}). \textit{Let $X$ be in} $\Cat$\textit{, with $b$, $c$ in $X$ such that}

\begin{itemize}
\item[\textit{a)}] \textit{$b\times c$ is in $X$.}
\item[\textit{b)}] \textit{$\forall$ $x$ in $X$, $b(x)$ ($\deq$} $\Hom_X(x,b)$\textit{) and $c(x)$ are $\neq\emptyset$, i.e. $X/b \longrightarrow X$ and $X/c \longrightarrow X$ are surjective on objects.}
\end{itemize}

\textit{Then $X$ is $\mathbf{W}_\infty$-aspheric.} (\footnote{False, cf. below. $X$ is 1-connected, but not necessarily 2-connected.})

The fact that $b$ is a greatest element of $\Ob X$ already ensures 0-connectedness.

For 1-connectedness, let us prove that in the fundamental groupoid $\Pi(X)$, $a \deq b\times c$ becomes a final element.
\[
\begin{tikzcd}
b\times c \arrow[r] \arrow[d] & b \\
c &
\end{tikzcd}
\]

For this, let us first prove that $\forall$ $x$ in $X$, all the elements of $a(x) = b(x)\times c(x)$ are equal in $\Fl\Pi(X)$. Indeed, since $b\times c \longrightarrow b$ is \texttt{[an]} isomorphism, hence \texttt{[a]} monomorphism in $\Pi(X)$, this already shows that if $\beta,\beta'\in b(x)$, $\gamma,\gamma'\in c(x)$, whence $(\beta,\gamma)$, $(\beta,\gamma')$, $(\beta',\gamma')\in b(x)\times c(x)$, we have $(\beta,\gamma) = (\beta,\gamma')$ in $\Pi(X)$ (since $\pr_1(\beta,\gamma) = \pr_1(\beta,\gamma')$, and likewise (using that $\pr_2$ is \texttt{[an]} isomorphism, hence \texttt{[a]} monomorphism) $(\beta,\gamma') = (\beta',\gamma')$). Thus all the $(\beta,\gamma)\in b(x)\times c(x) = a(x)$ are equal in $\Fl\Pi(X)$. To finish proving that $a$ is a final object of $\Pi(X)$, it suffices to see that every arrow \textit{in} $\Pi(X)$ from an $x\in\Ob X$ to $a$, comes from an arrow $x\longrightarrow a$ in $X$. But an arrow in $\Pi(X)$ is obtained by a path in $X$
\[
\begin{tikzcd}
x = x_0 \arrow[r, "u_1"] \arrow[dr, "\alpha_0"] & x_1 \arrow[d, "\alpha_1"'] & x_2 \arrow[l, "u_2"'] \arrow[r] \arrow[dl, "\alpha_2"'] & {} & \cdots & {} & x_{2N} = a \arrow[l, "u_{2N}"'] \\
& a. & & & & &
\end{tikzcd}
\]

Let us take an arrow $\alpha_0 : x\longrightarrow a_0$ in $X$, I say that the % [sic] a_0 -> a, the arrows α_i go to a (no a_0 is defined)
\msp{page 2}
homotopy class of the path $c$ (i.e. its image in $\Fl(\Pi(X))$ is equal to that of $u_0$. % [sic] u_0 -> \alpha_0 (?), no u_0 is defined; also the parenthesis opened at "(i.e." is not closed
To see this, let us choose for every $i\in[0,2N]$ an arrow $\alpha_i : x_i\longrightarrow a$. (It exists thanks to hypothesis b)). I say that for each $i\in[0,2N]$, the triangle
\printedpage{2}
% head: A. Grothendieck : Dérivateurs, Ch. XII, Édition des Universités de Montpellier II et Paris VII
% foot: Version 20 juin 2003
\[
\begin{tikzcd}[column sep=small]
x_0 \arrow[rr, "c_i"] \arrow[dr, "\alpha_0"] & & x_i \arrow[dl, "\alpha_i"'] \\
& a &
\end{tikzcd}
\]
in $\Pi(X)$ is commutative, where $c_i$ denotes the homotopy class of the path from $x_0$ to $x_i$ defined by the $u_i$ ($i\in[1,i]$) \texttt{[rather} ($i\in[1,2N]$)\texttt{]}. We do this by induction on $i$, the assertion is trivial if $i = 0$ ($c_0$ represents the identity map). The diagram % [sic] identitique -> identique, non-existent word
\[
\begin{tikzcd}[column sep=small]
x_0 \arrow[rr, "c_i"] \arrow[dr, "\alpha_0"] \arrow[rrrr, bend left=32, "c_{i+1}"] & & x_i \arrow[dl, "\alpha_i"'] \arrow[rr, "u_{i+1}"] & & x_{i+1} \arrow[dlll, "\alpha_{i+1}"'] \\ % verified: solid curved arrow; bend increased for label clearance
& a & & &
\end{tikzcd}
\]
shows that it suffices to prove that for $i\in[0,2N-1]$ the triangle
\[
\begin{tikzcd}[column sep=small]
x_i \arrow[rr, "u_{i+1}"] \arrow[dr, "\alpha_i"'] & & x_{i+1} \arrow[dl, "\alpha_{i+1}"] \\
& a &
\end{tikzcd}
\]
is commutative in $\Pi(X)$. (\textbf{NB} $u_{i+1}$ is in the direct or retrograde direction, according as % [sic] directe -> direct, agreement with "sens"
$i$ is even or odd.) But this follows from the fact that the image of $\Hom_X(x_i,a)$ (resp. $\Hom_X(x_{i+1},a)$) % [sic] paire ou impaire -> pair ou impair, agreement with "i"
\msp{page 3}
in $\Fl\Pi(X)$ is reduced to a single element.

Thus we find that $X$ \textit{is} 1-\textit{connected.} So there is hope! It remains to prove
\[
H^i(X,M) = 0, \qquad i\geq 2, \quad M \text{ constant coefficient.}
\]

I have looked at the case when we have this: $b = c$, $\Ob X = \{b, b\times b\}$, $\Hom_X(b\times b)$ \texttt{[rather} $\Hom(b\times b, b)$\texttt{]} $= \{$ disjoint union of the sets $u\circ\pr_1$ and $u\circ\pr_2$, with $u\in\Hom_X(b,b) = G$ an arbitrary monoid $\}$. Then $X^\circ$ is fibered, via $I\longmapsto b^I$, over the full subcategory $S$ of $\Ens$ consisting of $\{0\}$ and $\{0,1\}$, and thus $X$ is cofibered (in particular proper) over $S^\circ \deq T$. (NB $S^\circ = T$ has an initial object, but not a final object. Its cohomology with values in an arbitrary presheaf is not always trivial, cf. below.)

We find \texttt{[``We find'' crossed out]}
\[
H^n(X,M) \Longleftarrow H^p(T, \underbrace{R^q f_*(M_X)}_{= \mathcal{H}^q}), \qquad \text{where } f : X \longrightarrow S^\circ
\]
\texttt{[}$n = p+q$\texttt{]}. I admit here (provisionally) this: % [sic] J'admet -> J'admets, verb ending
\[
\begin{aligned}
H^p(T,F) &= 0 \ , &&\text{for } p\geq 2, \text{ any } F\in T^\wedge \ ,\\
H^0(T,F) &= \Ker\delta \ , &&\\
H^2(T,F) &= \Coker\delta \ , && % [sic] H^2(T,F) -> H^1(T,F), the first line gives H^p(T,F) = 0 for p >= 2
\end{aligned}
\]
\printedpage{3}
% head: A. Grothendieck : Dérivateurs, Ch. XII, Édition des Universités de Montpellier II et Paris VII
% foot: Version 20 juin 2003
where $\delta : F_0 \longrightarrow F_1$ is defined as the difference $\delta = p_1^* - p_2^*$, where $p_1$, $p_2$ are the two projections of $\{0,1\}\simeq\{0\}^2$ in $T = S^\circ$. Moreover we have
\[
(\mathcal{H}^q)_0 = H^q(G,M), \qquad (\mathcal{H}^q)_1 = H^q(G\times G, M),
\]
and
\[
p_1^*,\ p_2^*\ :\ H^q(G,M) \rightrightarrows H^q(G{\times}G, M) \qquad \text{are the transposes of } \pr_1,\ \pr_2 : G\times G\longrightarrow G.
\]

\msp{page 4}
We readily check that
\[
H^0(T,\mathcal{H}^q) = 0 \qquad \text{if } q\geq 1,
\]
and there remains
\[
H^n(X,M) \simeq H^1(S,\mathcal{H}^{n-1}) = \Coker(\delta : H^{n-1}(G) \longrightarrow H^{n-1}(G\times G)). % [sic] H^1(S, ...) -> H^1(T, ...), the cohomology is taken on T = S°
\]
This gives, for $n = 2$,
\[
\begin{aligned}
H^2(X,M) &\simeq \Coker\Big(\delta : \Hom_\mon(G,M) \longrightarrow \underbrace{\mathit{Hom}_\mon(G\times G, M)}_{\Hom_\mon(G,M)\times\Hom_\mon(G,M)} \ \Big)\\
&\simeq \Coker\Big(u\longmapsto (u,-u) : H \longrightarrow H\times H\Big),
\end{aligned}
\]
where $H$ is the abelian group $\Hom_\mon(G,M)$.

Whence
\[
H^2(X,M) \ \simeq\ H = \Hom_\mon(G,M).
\]
Thus
\begin{quote}
$X$ is 2-connected if and only if the enveloping \textit{abelian} group of the monoid $G$ is zero.
\end{quote}

Thus, if $G$ is an abelian group $\neq 0$, $X$ \textit{is not} 2-connected.

\textbf{Corollary.} $\Psi = {}$\begin{tikzcd}[ampersand replacement=\&, row sep=tiny, column sep=small, baseline=(A.base)]
\& b_0 \\
|[alias=A]| a_0 \arrow[ur] \arrow[dr] \& \\
\& c_0
\end{tikzcd}
\textit{is} not \textit{totally} $\mathbf{W}_\infty$\textit{-aspheric; i.e. if we have a diagram} \begin{tikzcd}[ampersand replacement=\&, row sep=tiny, column sep=small, baseline=(B.base)]
\& b \\
|[alias=B]| a \arrow[ur] \arrow[dr] \& \\
\& c
\end{tikzcd}
\textit{in a category} $X\in\Ob\Cat$\textit{, such that its amalgamated sum in $X^\wedge$ is equal to $e_X$, it does not follow that $X$ is $\mathbf{W}_\infty$-aspheric.} (But perhaps will be willing to be 1-connected?)

\printedpage{4}
% head: A. Grothendieck : Dérivateurs, Ch. XII, Édition des Universités de Montpellier II et Paris VII
% foot: Version 20 juin 2003

\msp{page 5}
\textbf{Question.} \textit{Are there other totally} $\mathbf{W}_\infty$\textit{-aspheric objects of} $\Cat$\textit{, than the categories having a final object?} I now suspect not -- hence that one cannot improve substantially on the already known notion of an open Dwyer immersion.

In the diagram
\[
\begin{tikzcd}
U \arrow[r, hook, "i"] \arrow[d, "p"'] \arrow[dr, phantom, "\substack{\text{cocart.}\\ \text{and cart.}}"] & X \arrow[d, "q"'] & T \arrow[l, hook', dotted, "j"'] \arrow[d, dotted, "\text{iso}"] \\
U' \arrow[r, hook, "i'"] & X' & T', \arrow[l, hook', dotted, "j'"']
\end{tikzcd}
\tag*{$(*)$}
\]
let us take
\[
\begin{aligned}
U = \Psi, \quad & T = \{t\} = e \text{ final category,}\\
& t \text{ final object of } X,\\
& \text{i.e. } H_t = (u\longmapsto \Hom_X(u,t)) = e_U \in \Ob U^\wedge \ \mbox{(\footnotemark)}.
\end{aligned}
\]
\footnotetext{i.e. $X$ obtained from $U = \Psi$ by adding to it a final object $t$ -- thus $X\simeq\boldsymbol{\Delta}^1\times\boldsymbol{\Delta}^1$,
\[
\begin{tikzcd}[ampersand replacement=\&,
execute at end picture={
\draw[rounded corners] ([xshift=-10pt,yshift=10pt]A.center) -- ([xshift=-10pt,yshift=-10pt]C.center) -- ([xshift=10pt,yshift=-10pt]D.center) -- ([xshift=10pt,yshift=3pt]D.center) -- ([xshift=3pt,yshift=3pt]C.center) -- ([xshift=3pt,yshift=10pt]A.center) -- cycle;
\node (P) at ([xshift=-35pt,yshift=-22pt]C.center) {$\Psi$};
\draw[-stealth] (P) -- ([xshift=-11pt,yshift=-4pt]C.center);
}]
|[alias=A]| {} \arrow[r] \& {} \\
|[alias=C]| {} \arrow[u] \arrow[r] \& |[alias=D]| {} \arrow[u]
\end{tikzcd}
\ .
\]}
Let us take for $U'$ the category spoken of in the hypothetical (and false) ``proposition'' p. 1 with $b = c$, more precisely a category with two objects $b$ and $b\times b$, a group (or even a monoid) $G$ of automorphisms \texttt{[(or even of endomorphisms)]} of $b$, and the other arrows those which follow from these, requiring that $b\times b$ really be the cartesian square of $b$. We find a cofibered category (hence proper) over the category $T$ opposite to that of finite sets of the form $\{0\}\cup\{0,1\}$ (two objects) with fibers $B_G$ and $B_{G\times G}$ respectively ($B_\pi =$ category with a single object, and endomorphism monoid $\pi$). The functor
\[
U = \Psi \longrightarrow U'
\]
is the one given by the diagram

\msp{page 6}
\[
\begin{tikzcd}[row sep=small]
& b \\
b\times b \arrow[ur, "\pr_1"] \arrow[dr, "\pr_2"'] & \\
& b
\end{tikzcd}
\]

\end{document}
